\documentclass[times,twocolumn,final]{elsarticle}


%% Stylefile to load CAG template
\usepackage{cag}
\usepackage{framed,multirow}
\usepackage{array}
%% The amssymb package provides various useful mathematical symbols
\usepackage{amssymb}
\usepackage{latexsym}
\usepackage{caption}
\usepackage{subcaption}
\usepackage{amsmath}

% Following three lines are needed for this document.
% If you are not loading colors or url, then these are
% not required.
\usepackage{url}
%\usepackage{xcolor}
\definecolor{newcolor}{rgb}{.8,.349,.1}

\newcommand{\ergun}[1]{{\color{red} #1}}

\usepackage{hyperref}

\usepackage[switch,pagewise]{lineno} %Required by command \linenumbers below

\usepackage{tikz}
\usetikzlibrary{calc,fit,positioning}

\journal{Computers \& Graphics}

\begin{document}

% \verso{T. Yildiz and E. Akleman}
\verso{Anonymous}

\begin{frontmatter}

\title{Triply Periodic Mesh-Based Linked-Structures: Design and Construction of a Family of Periodic Links Based on Bravais Lattices and Corresponding Wigner-Seitz Cells}

% \author[1]{Tolga Yildiz}

% \author[2]{Ergun Akleman \corref{cor1}}
% \cortext[cor1]{Corresponding author: 
%   Tel.: +1-979-739-0094; }
% \emailauthor{ergun.akleman@gmail.com}{Ergun Akleman}
%\ead{example@email.com}

% \address[1]{Department of Computer Science and Engineering, Texas A\&M University, College Station, TX, 77840, USA}
% \address[2]{Visual Computing and Computational Media Section, School of Performance, Visualization and Fine Arts, Joint with Department of Computer Science and Engineering, Texas A\&M University, College Station, TX, 77840, USA}


%\received{1 February 2017}
\received{\today}
%%%% Do not use the below for submitted manuscripts
%\finalform{28 March 2017}
%\accepted{2 April 2017}
%\availableonline{15 May 2017}
%\communicated{S. Sarkar}

\begin{abstract}
In this paper, we present a constructive framework for designing a broad family of \textit{triply periodic linked-structures} derived from polygonal meshes embedded in three dimensions. A \textit{linked-structure} consists of closed 3D curves—called \textit{face-cycles}—which are derived by transforming the faces of a 2-manifold, non-manifold, or 3-manifold mesh into unknotted loops represented as quadratic B-spline curves. These face-cycles are interconnected using \textit{edge-cycles}, defined as circular connectors centered on mesh edges and oriented orthogonally to the edge direction.


A key insight of our approach is the distinction between \textit{fully linked-structures}, where all mesh edges are bundled, and more general \textit{\ergun{partially-linked structures with single-link property}}, which can be obtained by labeling and bundling only a subset of edges. This selective edge labeling significantly expands the design space, enabling a wide range of topologically distinct links to be constructed from the same underlying mesh. 

We focus on meshes generated by \textit{Bravais lattices} and their corresponding \textit{Wigner--Seitz cells}, which enable the creation of \textit{triply periodic structures} that repeat in three linearly independent directions. These periodic configurations can be compactly encoded and manipulated using \textit{cubical arrays}, providing a formal framework for pattern generation and topological transformation.

Our implementation demonstrates the feasibility of the proposed method through a custom software pipeline integrated into Blender. The resulting linked-structures are visually simple, topologically robust, and physically realizable, with potential applications in materials science, architecture, and visual computing.
\end{abstract}


\begin{keyword}
%% MSC codes here, in the form: \MSC code \sep code
%% or \MSC[2008] code \sep code (2000 is the default)
%\MSC 41A05\sep 41A10\sep 65D05\sep 65D17
%% Keywords
\KWD Links \sep 3-Manifold Meshes \sep Triply Periodic Structures \sep Shape Modeling and Design 
\end{keyword}

\end{frontmatter}

\begin{figure*}[htpb]
    \centering
    % \includegraphics[width=0.90\textwidth]{images/teaser}
    \begin{tikzpicture}[node distance=0mm]
    \node (c1) [inner sep=0mm]             {\includegraphics[width=0.17\linewidth]{images/ws/CP.jpg}};
    \node (c2) [inner sep=0mm, right= of c1] {\includegraphics[width=0.17\linewidth]{images/ws/HP.jpg}};
    \node (c3) [inner sep=0mm, right= of c2] {\includegraphics[width=0.17\linewidth]{images/ws/CF.jpg}};
    \node (c4) [inner sep=0mm, right= of c3] {\includegraphics[width=0.17\linewidth]{images/ws/OF.jpg}};
    \node (c5) [inner sep=0mm, right= of c4] {\includegraphics[width=0.17\linewidth]{images/ws/CI.jpg}};
    % \node (a1) [inner sep=0mm, below= 2mm of c1] {\includegraphics[width=0.17\linewidth]{images/ws/CP_WS.jpg}};
    % \node (a2) [inner sep=0mm, right= of a1] {\includegraphics[width=0.17\linewidth]{images/ws/HP_WS.jpg}};
    % \node (a3) [inner sep=0mm, right= of a2] {\includegraphics[width=0.17\linewidth]{images/ws/CF_WS.jpg}};
    % \node (a4) [inner sep=0mm, right= of a3] {\includegraphics[width=0.17\linewidth]{images/ws/OF_WS.jpg}};
    % \node (a5) [inner sep=0mm, right= of a4] {\includegraphics[width=0.17\linewidth]{images/ws/CI_WS.jpg}};
    \node (b1) [inner sep=0mm, below= 2mm of c1] {\includegraphics[width=0.17\linewidth]{images/ws/CP1.jpg}};
    \node (b2) [inner sep=0mm, right= of b1] {\includegraphics[width=0.17\linewidth]{images/ws/HP1.jpg}};
    \node (b3) [inner sep=0mm, right= of b2] {\includegraphics[width=0.17\linewidth]{images/ws/CF1.jpg}};
    \node (b4) [inner sep=0mm, right= of b3] {\includegraphics[width=0.17\linewidth]{images/ws/OF1.jpg}};
    \node (b5) [inner sep=0mm, right= of b4] {\includegraphics[width=0.17\linewidth]{images/ws/CI1.jpg}};
    \node (d1) [inner sep=0mm, below= 2mm of b1] {\includegraphics[width=0.17\linewidth]{images/ws_link_base/cp_colored.jpg}};
    \node (d2) [inner sep=0mm, right= of d1] {\includegraphics[width=0.17\linewidth]{images/ws_link_base/hp_colored.jpg}};
    \node (d3) [inner sep=0mm, right= of d2] {\includegraphics[width=0.17\linewidth]{images/ws_link_base/cf_colored.jpg}};
    \node (d4) [inner sep=0mm, right= of d3] {\includegraphics[width=0.17\linewidth]{images/ws_link_base/of_colored.jpg}};
    \node (d5) [inner sep=0mm, right= of d4] {\includegraphics[width=0.17\linewidth]{images/ws_link_base/ci_colored.jpg}};
    \node (e1) [inner sep=0mm, below= 2mm of d1] {\includegraphics[width=0.17\linewidth]{images/ws_link_base/cp_colored_unique.jpg}};
    \node (e2) [inner sep=0mm, right= of e1] {\includegraphics[width=0.17\linewidth]{images/ws_link_base/hp_colored_unique.jpg}};
    \node (e3) [inner sep=0mm, right= of e2] {\includegraphics[width=0.17\linewidth]{images/ws_link_base/cf_colored_unique.jpg}};
    \node (e4) [inner sep=0mm, right= of e3] {\includegraphics[width=0.17\linewidth]{images/ws_link_base/of_colored_unique.jpg}};
    \node (e5) [inner sep=0mm, right= of e4] {\includegraphics[width=0.17\linewidth]{images/ws_link_base/ci_colored_unique.jpg}};
    \node (f1) [inner sep=0mm, below= 2mm of e1] {\includegraphics[width=0.17\linewidth]{images/ws_linked/cp_linked.jpg}};
    \node (f2) [inner sep=0mm, right= of f1] {\includegraphics[width=0.17\linewidth]{images/ws_linked/hp_linked.jpg}};
    \node (f3) [inner sep=0mm, right= of f2] {\includegraphics[width=0.17\linewidth]{images/ws_linked/cf_linked.jpg}};
    \node (f4) [inner sep=0mm, right= of f3] {\includegraphics[width=0.17\linewidth]{images/ws_linked/of_linked.jpg}};
    \node (f5) [inner sep=0mm, right= of f4] {\includegraphics[width=0.17\linewidth]{images/ws_linked/ci_linked.jpg}};
    \node (g1) [inner sep=0mm, below= 2mm of f1] {\includegraphics[width=0.17\linewidth]{images/ws_linked/cp_linked1.jpg}};
    \node (g2) [inner sep=0mm, right= of g1] {\includegraphics[width=0.17\linewidth]{images/ws_linked/hp_linked1.jpg}};
    \node (g3) [inner sep=0mm, right= of g2] {\includegraphics[width=0.17\linewidth]{images/ws_linked/cf_linked1.jpg}};
    \node (g4) [inner sep=0mm, right= of g3] {\includegraphics[width=0.17\linewidth]{images/ws_linked/of_linked1.jpg}};
    \node (g5) [inner sep=0mm, right= of g4] {\includegraphics[width=0.17\linewidth]{images/ws_linked/ci_linked1.jpg}};

    
    \node[on grid, align=center, inner sep=1mm,draw, rounded corners=1mm,text height=18pt, minimum width=31mm, above=24mm of c1](c1_l){Primitive\\Cubic (cP)};
    \node[on grid, align=center, inner sep=1mm,draw, rounded corners=1mm,text height=18pt, minimum width=31mm, above=24mm of c2](c2_l){Primitive\\Hexagonal (hP)};
    \node[on grid, align=center, inner sep=1mm,draw, rounded corners=1mm,text height=18pt, minimum width=31mm, above=24mm of c3](c3_l){Primitive Rhombo-\\hedral (hR) $\alpha<90^\circ$};
    \node[on grid, align=center, inner sep=1mm,draw, rounded corners=1mm,text height=18pt, minimum width=31mm, above=24mm of c4](c4_l){Primitive\\Triclinic (aP)};
    \node[on grid, align=center, inner sep=1mm,draw, rounded corners=1mm,text height=18pt, minimum width=31mm, above=24mm of c5](c5_l){Primitive Rhombo-\\hedral (hR) $\alpha>90^\circ$};
    
    \path node (l1) [draw,align=center,on grid,rotate=90,rounded corners, left=22mm of c1, minimum width=31mm, minimum height=6mm]{LATTICE + WS};
    % \path node (l2) [draw,align=center,on grid,rotate=90,rounded corners, left=22mm of a1, minimum width=31mm, minimum height=6mm]{TESSELATION};
    \path node (l3) [draw,align=center,on grid,rotate=90,rounded corners, left=22mm of b1, minimum width=31mm, minimum height=6mm]{WS CELL};
    \path node (l4) [draw,align=center,on grid,rotate=90,rounded corners, left=22mm of d1, minimum width=31mm, minimum height=6mm]{LINK SETS};
    \path node (l5) [draw,align=center,on grid,rotate=90,rounded corners, left=22mm of e1, minimum width=31mm, minimum height=6mm]{UNIQUE LINKS};
    \path node (l6) [draw,align=center,on grid,rotate=90,rounded corners, left=22mm of f1, minimum width=31mm, minimum height=6mm]{FULLY-LINKED};
    \path node (l7) [draw,align=center,on grid,rotate=90,rounded corners, left=22mm of g1, minimum width=31mm, minimum height=6mm]{PARTIAL};

    % \node[draw, fit=(a1)(a5), rounded corners=1mm, inner sep=0mm]{};
    \node[draw, fit=(b1)(b5), rounded corners=1mm, inner sep=0mm]{};
    \node[draw, fit=(c1)(c5), rounded corners=1mm, inner sep=0mm]{};
    \node[draw, fit=(d1)(d5), rounded corners=1mm, inner sep=0mm]{};
    \node[draw, fit=(e1)(e5), rounded corners=1mm, inner sep=0mm]{};
    \node[draw, fit=(f1)(f5), rounded corners=1mm, inner sep=0mm]{};
    \node[draw, fit=(g1)(g5), rounded corners=1mm, inner sep=0mm]{};
    \end{tikzpicture}
\caption{Overview of five primitive Bravais lattices and their corresponding linked-structures.
Top row: Voronoi tessellations of lattice points showing Wigner–Seitz (WS) cells.
Second row: The five \ergun{combinatorially} distinct WS cell shapes derived from each lattice: left to right: cube, hexagonal prism, rhombic dodecahedron, elongated dodecahedron, truncated octahedron.
Third row: Linked-structures formed from individual WS cells using fully linked bundling.
Fourth row: Unique link segments \ergun{(shown in different colors)} identified per cell for minimal redundancy. \ergun{A brute-force algorithm is used to identify the unique links.}
Fifth row: The fully-linked configurations.
Bottom row: Partially-linked single-link configurations using a minimal set of edge-cycles.
 }
\label{fig_teaser}
\end{figure*}

%\linenumbers


\section{Introduction}

An important class of structures with broad applications in science and engineering are \textit{geometrically (or mechanically) interlocked architectures}. These architectures—such as mathematical links and chainmail-inspired designs—are gaining prominence in science and engineering due to their unusual blend of flexibility, strength, and customizable mechanical behavior \cite{carlucci2003polycatenation,zhou20253d,cheng20253d,zhang2025recognizing,tang2023chainmail}.

In chemistry, these structures are categorized under the general term \textit{mechanically interlocked molecular architectures}, which include \textit{catenanes}, \textit{rotaxanes}, \textit{molecular knots}, and \textit{links} (Figure~\ref{fig_MIS}). Recent advances in electrochemical switching mechanisms \cite{electrochem2024} and molecular handcuff designs \cite{handcuffs2022} have expanded their potential for molecular machinery and functional materials.

In materials science, these architectures serve as the foundation for metamaterials—engineered systems whose mechanical properties arise not just from material composition, but from their internal geometry. Chainmail-like lattices, for instance, are being used to produce lightweight, flexible, and impact-resistant materials suitable for protective equipment \cite{nature2016interlocked,iass2024parametric} and biomedical applications \cite{pmc2024deployable}. Recent developments in 3D-printed polycatenated architected materials (PAMs) demonstrate unprecedented energy dissipation capabilities through their semi-liquid structural dynamics \cite{popularmechanics2025}.

In fields like mechanical and aerospace engineering, interwoven mesh systems offer efficient stress distribution and energy absorption, making them well-suited for shock-absorbing parts or deployable structures \cite{pmc2024deployable}. Their ability to deform—bend, stretch, or compress—while retaining connectivity makes them ideal for use in soft robotics, prosthetics, and responsive surfaces \cite{arxiv2024chainmail}. DNA-based molecular chainmail structures show particular promise for nanoscale engineering applications \cite{arxiv2024chainmail}.

Despite these advances, a major challenge remains: the lack of systematic \textit{shape modeling} principles for creating such interlocked structures. Most current methods rely on heuristic or ad hoc strategies, emphasizing fabrication and physical testing over foundational design theory. Our work addresses this gap by introducing a mathematical framework grounded in discrete geometry and symmetry. \ergun{By leveraging a subset of 3D lattice symmetries, we propose a method for systematically classifying and designing these structures. This approach has the potential to uncover structural insights that might remain hidden in purely empirical or numerical analyses.}

\begin{figure}[htpb]
\centering
\includegraphics[width=0.99\columnwidth]{images/MIS}
    \caption{\it Examples of mechanically interlocked architectures. From left to right: A catenane (interlocked graph), a rotaxane (graph with a geometric barrier), a link (set of closed, interlinked curves embedded in 3D).}
\label{fig_MIS}
\end{figure}

%\subsection{Background on Linked-Structures}
\subsection{Basis and Rationale}\label{sec:basis}
\label{BasisAndRationale}

The geometrically (or mechanically) interlocked architectures can be mathematically described as \textit{embeddings of a set of disconnected graphs in three-dimensional Euclidean space}, $\mathbb{R}^3$. \ergun{In mathematical terminology, when a graph consists of a single simple cycle, any of its 3D embeddings is called a \textit{knot} \cite{neuwirth1979theory}.} If multiple disjoint cycles (i.e., multiple unlinked or interlinked circles) are embedded in $\mathbb{R}^3$, the structure is referred to as a \textit{link}. When the individual cycles are knotted, the structure is a \textit{knotted link}. This category includes \textit{woven objects}~\cite{akleman2009cyclic,akleman2011cyclic,xing2010single,akleman2020topologically,akleman2015extended} and \textit{knitted objects}~\cite{knittel2020modelling,wadekar2021geometric,niu2025geometric}. If all embedded cycles are \textit{unknots} (i.e., topologically equivalent to a circle), the resulting structure is referred to simply as a \textit{link}—examples include the \textit{Hopf link}, \textit{Solomon’s knot}, \textit{Whitehead link}, and the \textit{Borromean rings} (Figure~\ref{fig_famouslinks}).

\begin{figure}[htpb]
\centering
\includegraphics[width=0.99\columnwidth]{images/simplelinks}
    \caption{\it Examples of classical mathematical links. From left to right: Hopf link, Borromean rings, Whitehead link, and Solomon’s knot—all composed of unknotted but interlinked closed curves. }
\label{fig_famouslinks}
\end{figure}

While links are naturally a subset of knotted structures \ergun{that can be produced by twisting edges of manifold meshes ~\cite{akleman2020topologically,akleman2015extended}, these edge twisting methods do not guarantee} the generation of \textit{link-only} configurations without introducing knots.
As a motivating example, consider the Borromean rings. They can be produced by applying a uniform half-twist to all edges of a tetrahedron using graph rotation systems. However, without formal guidance, the resulting geometry may not resemble the traditional depiction of Borromean rings. The geometry could deviate significantly, making the structure difficult to recognize visually. Although such constructions are topologically valid, they often fail to yield \textit{geometrically simple and interpretable 3D forms}—such as the familiar arrangement of three ellipsoids in space. As a result, while suitable for topological exploration, these approaches often lack the geometric clarity and visual control needed for engineering, education, and artistic applications.

\subsection{Context \& Motivation}

Our objective is to establish a \textit{large and constructive design space} for the systematic generation of linked-structures.
\ergun{To achieve this objective we first introduce a general replacement system that converts polygonal meshes into linked structures composed of \textit{trivial knots}—that is, unknots that are visually recognizable. Figure~\ref{fig_teaser1b} illustrates our replacement system, which generates trivial links from any 2-manifold, non-manifold, or 3-manifold mesh with appropriately “labeled” edges, in a nutshell.} The method is based on representing mesh faces as trivial knots, or \textit{face-cycles}, using closed quadratic B-Spline curves derived from face boundaries. These face-cycles are bundled together using \textit{edge-cycles}, which are circles centered at the midpoint of labeled edges of given mesh and oriented perpendicular to the edge directions. The radius of each edge-cycle is selected to be large enough to fully encompass all corresponding face-cycles (Figure~\ref{fig_teaser1b}). 


\begin{figure}[htpb]
    \centering        
        \begin{subfigure}[t]{0.31\columnwidth}
    \includegraphics[width=0.99\linewidth]{images/link_algo/2man}
            \caption{\it The original mesh.}
        \label{fig_3D_Mesh0_0}
    \end{subfigure}
        \hfill
       \begin{subfigure}[t]{0.31\columnwidth}
    \includegraphics[width=0.99\linewidth]{images/link_algo/2man_linked}
            \caption{\it Related fully-linked structure.}
        \label{fig_3D_Mesh0_1}
    \end{subfigure}
        \hfill
         \begin{subfigure}[t]{0.31\columnwidth}
    \includegraphics[width=0.99\linewidth]{images/link_algo/2man_linked_partial}
            \caption{\it A \ergun{partially-linked} structure with single-link property.}
        \label{fig_3D_Mesh0_2}
    \end{subfigure}  
\caption{\it An example of converting polygonal meshes into linked-structures using face smoothing and edge-cycle bundling. In this case, we converted a regular dodecahedral mesh with convex planar faces into linked-structures with visually clear links. Note that Figure~\ref{fig_3D_Mesh0_2} remains geometrically interlocked, even though not all face-cycles are bundled. }
\label{fig_teaser1b}
\end{figure}

\ergun{A key strength of our approach is that any mesh with labeled edges uniquely defines a corresponding linked structure. This allows us to analyze whether the resulting linked structure is geometrically interlocked directly from the labeled mesh. For the simplest form of geometric interlocking, we focus on structures consisting of a single connected component, which inherently prevents disassembly into smaller, disconnected parts. We refer to this characteristic as the single-link property, and we call such linked structures single-links. Importantly, achieving single-links does not require linking all face-cycles, as demonstrated in Figure~\ref{fig_3D_Mesh0_2}. This property is valuable both mathematically and practically, as it supports the design of cohesive, globally connected structures while maintaining design flexibility.} 

\subsection{Visual Clarity, Planarity, and Periodicity with Triply-Periodic Wigner-Seitz Cells}

\ergun{Note that this approach is highly general and can be applied to a wide variety of linked structures. However, we realized that such generality is not always beneficial for the systematic construction of linked configurations. Through experimentation, we observed that it is more effective to begin with meshes composed of planar faces.}

It is important to recognize that visual crossings may occur depending on the viewer’s perspective, particularly when curves lie on non-planar surfaces. For example, circular curves traced on bilinear surfaces—such as hyperbolic paraboloids—can appear to intersect from certain viewpoints, even when they are topologically unknotted.

To reduce such visual ambiguity, we prioritize constructing meshes with \textit{planar faces} whenever feasible. Fortunately, this is often achievable; for instance, the faces of Voronoi cells are guaranteed to be planar when the sites are points and the metric is Euclidean. By converting the boundaries of these faces into \ergun{planar curves in 3D}, we ensure that the resulting links are visually clean and free of apparent crossings from all viewpoints.


\ergun{We are also interested in using these structures for architecturing materials in engineering. Periodicity plays a central role in the design and functionality of architected materials. Periodic structures offer predictable mechanical, thermal, and optical properties, enabling precise control over material behavior at multiple scales. The periodicity not only ensures uniform distribution of stress and deformation under external loads but also facilitates manufacturability through tiling and modular construction.

Moreover, periodicity allows the incorporation of topological features—such as single-link configurations—into the material’s internal geometry without compromising its structural regularity. By aligning the combinatorial design of links with lattice symmetry, we can create a versatile design space that can support the development of novel metamaterials with tunable stiffness, flexibility, and failure resistance. }


In this paper, we introduce a methodology for designing a specific class of periodic links, which we call \textit{triply periodic linked-structures}. These structures are constructed by transforming triply periodic Wigner-Seitz cells into link configurations (Figure~\ref{fig_teaser}). This provides a systematic framework for generating mechanically interlocked configurations with global connectivity.  An additional and important property of  Wigner-Seitz cells is that they are produced by Voronoi decomposition, and their faces are naturally planar.

\section{Related Work} 

We recently worked on designing a novel repeating link, leading to a method for filling space using Borromean links. However, this ad hoc approach lacked generalization, crucial for systematically creating new links and advancing engineering applications.

Several systematic methods exist for designing knots and links, using guiding shapes like graphs \cite{tait1877knots, tait1877links}, polyhedra \cite{roelofs2007rings,sequin2021polyhedral,lipschutz2022linkedknots,gailiunas2022tplinks}, extended graph rotation systems \cite{akleman2020topologically, akleman2015extended}, and CMM-pattern coverable meshes \cite{akleman2015block}. While effective in generating topologically diverse links and knots, these methods offer limited geometric control.

Recent advances in machine learning have introduced new computational approaches for knot recognition and generation. Transformer-based models now achieve almost 99\% accuracy in knot-type classification while being 4500× faster than traditional Alexander polynomial methods \cite{zhang2025recognizing}. Multiscale Gauss link integrals (mGLI) provide enhanced topological analysis of protein structures through curve segmentation and multiscale analysis \cite{shen2023knot}. These computational advances complement our geometric approach by enabling rapid verification of topological properties in generated structures.

The motivation for this work originates from the design of chainmail structures, a type of body armor historically worn by European knights during the Middle Ages \cite{grancsay1928notes}. Modern computational studies have quantified the relationship between module quantity and structural curvature in 3D chainmail systems \cite{afif2024influence}, while advances in 3D printing enable fabrication of quasi-rigid interlocking elements with remarkable flexibility-strength combinations \cite{tang2023chainmail}. Unlike traditional chainmail, which consists of interlinked trivial knots, modern chainmail structures can be formed using 1-complexes, such as the edges of a cube \cite{afif2024chain,wang2021structured,xu2023study}.

While chainmail structures inspire this study, the focus of this paper is exclusively on linked trivial knots, which have been a part of formal exploration of knots and links, which began in the 19th century with pioneering work by Lord Kelvin, Peter Guthrie Tait, and James Clerk Maxwell, who initiated efforts to classify and understand the structure of knots and links \cite{tait1877links,tait1877knots}. During the 20th century, knot theory developed into a major area within topology, with significant advances in the classification of knots \cite{conway1970enumeration,neuwirth1979theory}.

Today, knot theory extends beyond pure mathematics and finds diverse applications in fields such as biology and chemistry \cite{zhang1994dnatroctahedron,moestopo2023knots,fielden2017molecular,craik2001cystine,horner2016knot,forgan2011chemical}. It also continues to influence visual arts, sculptural design \cite{verhoeff2011chainlinkfence,holden2016braids,sequin2021polyhedral,lipschutz2022linkedknots}, and architectural structures \cite{kozlov2013structures,borhani2016emergence,borhani2016material,kozlov2018knots}.

Our method in this paper is essentially a replacement method that uses guide shapes to create knots and links \cite{roelofs2007rings,sequin2021polyhedral,gailiunas2022tplinks}. It is an extended version of the Medial Graph method introduced by Peter Guthrie Tait \cite{tait1877links,tait1877knots}, where each vertex of the medial graph is replaced by a crossing to form knots and links. However, a problem with medial graphs is that the same graph can correspond to different links, as the graph can be embedded in topologically different surfaces \cite{gross2001topological}. To resolve this ambiguity, graph rotation systems are employed \cite{gross2020enumerating,akleman1999guaranteeing}. Based on extended graph rotation systems, knots and links can be obtained as boundaries of faces of non-orientable surfaces immersed in 3D space \cite{akleman2020topologically,akleman2015extended}. The crossings are achieved by twisting the edges of orientable 2-manifold meshes. More complex knots and links can be constructed by twisting the edges of 3-manifold meshes \cite{akleman2015block}. This approach is used to create plain woven objects, single-cycle plain woven objects, and twill-woven objects \cite{akleman2009cyclic,xing2010single,akleman2011cyclic}.

Recent developments in high-quality mesh generation \cite{liu2024meshformer} have enabled new approaches for 3D structure creation using sparse-view reconstruction models with explicit 3D native structure guidance. More intricate structures can also be generated using Voronoi decompositions of space, with knotted curves and links serving as Voronoi sites \cite{krishnamurthy2020bi,krishnamurthy2021geometrically,yildiz2023modular,yildiz2024volumetric}.

Despite their power, these methods have a significant limitation: they cannot produce geometrically simple structures. In this paper, we present a method to obtain geometrically simple linked-structures. These simply linked-structures resemble soft cells as they consist solely of curves that are tangent to each other \cite{domokos2024soft}. Our method builds upon a previous approach that can generate fractal soft cells \cite{akleman2004smooth}. The process involves converting face boundaries into quadratic B-Spline curves, which are designed to ensure tangency to each edge of the face at the midpoint of the edges. This property guarantees the characteristic soft cell appearance.



\section{Periodic Meshes with Labeled Edges: An Overview}

\ergun{This section provides a conceptual overview of periodic meshes with labeled edges and their transformation into linked structures. To ease readers into the topic of triply periodic constructions, we begin with the simpler case of doubly periodic 2-manifold meshes. We then present concrete examples of such meshes and demonstrate how their edge labeling leads to corresponding linked structures.

The ideas introduced here—including the formulation of periodic meshes, labeling strategies, and the combinatorial conversion into links—are, to the best of our knowledge, original contributions and are being presented for the first time. Given the novelty and potential generality of this framework, we aim in this section to provide a clear and structured exposition that helps readers understand the core concepts and appreciate the breadth of the design space it enables.}

\begin{figure}[htpb]
\centering
% \includegraphics[width=0.30\columnwidth]{images/2d_wigner/square.jpg}
% \includegraphics[width=0.30\columnwidth]{images/2d_wigner/hex.jpg}
% \includegraphics[width=0.30\columnwidth]{images/2d_wigner/square1.jpg}
\includegraphics[width=0.30\columnwidth]{images/2DWS0}
\includegraphics[width=0.30\columnwidth]{images/2DWS1}
\includegraphics[width=0.30\columnwidth]{images/2DWS2}
    \caption{\it Examples of 2D Bravais lattices (black points) and their corresponding Wigner–Seitz cells (pink polygons). These generate doubly periodic planar 2-manifold meshes that tile the Euclidean plane without boundaries—useful for understanding the triply periodic 3D cases.}
\label{fig_2DWS}
\end{figure}

\subsection{Doubly-Periodic 2-Manifold Meshes}

Figure~\ref{fig_2DWS} illustrates several examples of two-dimensional Wigner–Seitz cells, obtained through Voronoi decomposition of 2D Euclidean space using 2D Bravais lattices as Voronoi sites. We observe that many traditional chainmail structures used in armor can be reproduced from these doubly periodic meshes via the replacement scheme shown in Figures~~\ref{fig_teaser1b}, \ref{fig_teaser1a}, and ~\ref{fig_teaser1c}. In this sense, the approach presented in this paper can be seen as a formalization and generalization of 2D chainmail into three-dimensional linked-structures.




\begin{figure}[htpb]
    \centering
        \begin{subfigure}[t]{0.24\columnwidth}
  \includegraphics[width=1.0\textwidth]{images/2D/00}
        \caption{\it No edge labeled.  }
        \label{fig_2D/00}
    \end{subfigure}
    \hfill
            \begin{subfigure}[t]{0.24\columnwidth}
  \includegraphics[width=1.0\textwidth]{images/2D/01}
        \caption{\it One edge labeled.  }
        \label{fig_2D/01}
    \end{subfigure}
        \hfill
    \begin{subfigure}[t]{0.24\columnwidth}
  \includegraphics[width=1.0\textwidth]{images/2D/10}
        \caption{\it One edge labeled.  }
        \label{fig_2D/10}
    \end{subfigure}
                \begin{subfigure}[t]{0.24\columnwidth}
  \includegraphics[width=1.0\textwidth]{images/2D/11}
        \caption{\it Two edges.  }
        \label{fig_2D/11}
    \end{subfigure}
\caption{\it Four edge-labeling configurations for a square grid based on a 2D Bravais lattice.
Green edges are labeled (linked); red edges are unlabeled. Thin black edges indicate labels inferred from adjacent square tiles, ensuring continuity across cells. }
\label{fig_teaser2}
\end{figure}



Building on this formalization, we present two methods for generating periodic 3-manifold meshes that can be converted into triply periodic linked-structures within the 3D Bravais lattice framework. The first method leverages 3D Wigner–Seitz cells, which provide a natural and rigorous representation of repeating units capable of uniquely defining link patterns in volumetric space.  To create these patterns, we selectively label the edges of the congruent polyhedra. 


The second method generates random planar meshes that tile space through repetition along three linearly independent directions. These meshes are constructed with topological constraints that ensure continuity and compatibility across repeating units, thereby achieving seamless periodicity. Note that a simple decomposition of a space-filling polyhedron is not sufficient to produce a triply periodic mesh. For example, the decomposed square and cube shown in Figures~\ref{fig_teaser1a} and~\ref{fig_teaser1c} are not doubly or triply periodic random meshes, as their opposite boundaries do not align. The method for generating random triply periodic meshes is described in Section~\ref{sec_TP3MMwPF}.

\begin{figure}[htpb]
    \centering
        \begin{subfigure}[t]{0.24\columnwidth}
  \includegraphics[width=1.0\textwidth]{images/2D/000}
        \caption{\it No edge labeled.  }
        \label{fig_2D/000}
    \end{subfigure}
    \hfill
            \begin{subfigure}[t]{0.24\columnwidth}
  \includegraphics[width=1.0\textwidth]{images/2D/001}
        \caption{\it One edge labeled.  }
        \label{fig_2D/001}
    \end{subfigure}
        \hfill
    \begin{subfigure}[t]{0.24\columnwidth}
  \includegraphics[width=1.0\textwidth]{images/2D/010}
        \caption{\it One edge labeled.  }
        \label{fig_2D/010}
    \end{subfigure}
                \begin{subfigure}[t]{0.24\columnwidth}
  \includegraphics[width=1.0\textwidth]{images/2D/011}
        \caption{\it Two edges.  }
        \label{fig_2D/011}
    \end{subfigure}
    \hfill
    \begin{subfigure}[t]{0.24\columnwidth}
  \includegraphics[width=1.0\textwidth]{images/2D/100}
        \caption{\it One edge labeled.  }
        \label{fig_2D/100}
    \end{subfigure}
    \hfill
    \begin{subfigure}[t]{0.24\columnwidth}
  \includegraphics[width=1.0\textwidth]{images/2D/101}
        \caption{\it Two edges.  }
        \label{fig_2D/101}
    \end{subfigure}
        \hfill
    \begin{subfigure}[t]{0.24\columnwidth}
  \includegraphics[width=1.0\textwidth]{images/2D/110}
        \caption{\it Two edges.  }
        \label{fig_2D/110}
    \end{subfigure}
        \begin{subfigure}[t]{0.24\columnwidth}
  \includegraphics[width=1.0\textwidth]{images/2D/111}
        \caption{\it Three edges.  }
        \label{fig_2D/111}
    \end{subfigure}
\caption{\it All eight labeling cases for hexagonal grids derived from 2D Bravais lattices.
As in Figure 6, red = unlabeled, green = labeled.
Labels on shared edges are inferred from neighboring hexagonal cells. These configurations form the basis of doubly periodic hexagonal meshes. }
\label{fig_teaser4}
\end{figure}

\subsection{Edge Labeling}

We also demonstrate the range of design options for creating repeating patterns that can be further expanded \ergun{with periodic structures}. In particular, \ergun{it is possible to use a pattern} for assigning a label to every edge of the original mesh to obtain geometrically interlocked \textit{partially-linked structures}. \ergun{Such partially-linked structures with single-link property} can often be achieved through many different configurations that use only a selected subset of edges, as illustrated in Figures~\ref{fig_teaser3} and~\ref{fig_teaser5}.

Moreover, since the underlying lattice structure is Bravais and each polyhedron can be chosen to be congruent, edge-cycles can be systematically assigned using additional patterning strategies. For lattice-based periodic meshes, we observe that it is possible to assign labels to each lattice point that uniquely determine which edges should be linked (see Figures~\ref{fig_teaser2} and ~\ref{fig_teaser4}).




\begin{figure}[htpb]
    \centering
        \begin{subfigure}[t]{0.48\columnwidth}
  \includegraphics[width=1.0\textwidth]{images/2D/mesh1}
        \caption{\it A labeled square grid defined by doubly periodic $2x2$ region.}
        \label{fig_2D/mesh1}
    \end{subfigure}
    \hfill
            \begin{subfigure}[t]{0.48\columnwidth}
  \includegraphics[width=1.0\textwidth]{images/2D/link1}
        \caption{\it A doubly periodic link that corresponds the mesh in \ref{fig_2D/mesh1}.  }
        \label{fig_2D/link1}
    \end{subfigure}
\caption{\it A \ergun{partially-linked structure with single-link property} on a square grid.
(a) A doubly periodic 2×2 labeling pattern.
(b) The resulting link, where not all edges are bundled, yet the structure remains connected as a single topological link.  }
\label{fig_teaser3}
\end{figure}

\ergun{Whether the resulting partially linked structure is geometrically interlocked (i.e., possesses the single-link property) can be determined by constructing a dual mesh in which each edge is explicitly labeled as either \textit{linked} or \textit{unlinked}, based on its status in the original mesh.
This selective linking introduces additional degrees of freedom, significantly expanding the overall \textit{design space} and enabling a richer variety of topological and geometric configurations.}




\begin{figure}[htpb]
    \centering
        \begin{subfigure}[t]{0.48\columnwidth}
  \includegraphics[width=1.0\textwidth]{images/2D/mesh2}
        \caption{\it A labeled hexagonal grid defined by doubly periodic $2x2$ region.}
        \label{fig_2D/mesh2}
    \end{subfigure}
    \hfill
            \begin{subfigure}[t]{0.48\columnwidth}
  \includegraphics[width=1.0\textwidth]{images/2D/link2}
        \caption{\it A doubly periodic link that corresponds the mesh in \ref{fig_2D/mesh2}.  }
        \label{fig_2D/link2}
    \end{subfigure}
\caption{\it A \ergun{partially-linked structure with single-link property} on a hexagonal Bravais grid.
(a) Labeling pattern for a 2×2 unit cell on a hexagonal mesh.
(b) Resulting link, with connectivity preserved across cells via consistent labeling strategy. The underlying Bravais symmetry ensures periodicity. }
\label{fig_teaser5}
\end{figure}

\subsection{Labeling Strategies}

A key advantage of using the Bravais lattice framework is that edge connections can be labeled systematically using cubical arrays similar to those in traditional fabric patterns (see examples for doubly periodic 2-manifold meshes in Figures~\ref{fig_teaser3} and ~\ref{fig_teaser5}). These arrays uniquely encode link patterns and can be further manipulated through operations involving 3D tensors defined within $k$-adic zero algebra. This algebraic representation enables structured control over the design and transformation of linked mesh patterns.

\section{Preliminaries and Theoretical Framework}
\label{sec_preliminaries}


This section lays the mathematical foundation for our approach by formally introducing the mesh-to-link conversion process and the role of lattice-based structures in generating periodic links. We begin demonstrate 
how any polygonal mesh can be converted into a linked-structure using face-cycles and edge-cycles. We then narrow our focus to triply periodic 3-manifold meshes derived from Bravais lattices and their corresponding Wigner–Seitz cells, which provide a rigorous basis for encoding repetition and symmetry in three-dimensional space. Finally, we define volumetric arrays and label-based configurations that allow these periodic meshes to be transformed systematically into triply periodic linked-structures.

We refer to a structure as \textit{triply (or doubly) periodic} if it exhibits translational symmetry along three (or two) linearly independent directions, i.e., it is periodic with respect to a Bravais lattice in $\mathbb{R}^3$. For 3D linked-structures, it is desired to have such translational symmetry with repeating patterns, similar to those found in woven or knitted materials, which is particularly valuable for the design of architectured materials. To achieve this, we construct links based on cell-transitive honeycombs, which are generated through translations of a single polyhedron—a special class of space-filling congruent shapes.

To create these congruent polyhedral cells, we utilize translational symmetries defined by the 3D Bravais lattices. Specifically, we construct three-dimensional cell-transitive honeycombs using Wigner–Seitz cells, which are derived via Voronoi decomposition of the 3D Euclidean space with Bravais lattice positions serving as Voronoi sites.


\begin{figure}[htpb]
    \centering        
        \begin{subfigure}[t]{0.31\columnwidth}
    \includegraphics[width=0.99\linewidth]{images/link_algo/2manwb}
            \caption{\it The original mesh.}
        \label{fig_2D_Mesh0_0}
    \end{subfigure}
        \hfill
       \begin{subfigure}[t]{0.31\columnwidth}
    \includegraphics[width=0.99\linewidth]{images/link_algo/2manwb_linked}
            \caption{\it Related  fully-linked structure.}
        \label{fig_2D_Mesh0_1}
    \end{subfigure}
        \hfill
         \begin{subfigure}[t]{0.31\columnwidth}
    \includegraphics[width=0.99\linewidth]{images/link_algo/2manwb_linked_partial}
            \caption{\it A partially-linked structure.}
        \label{fig_2D_Mesh0_2}
    \end{subfigure}
\caption{\it An example of a random 2-manifold mesh with boundary and its conversion into a linked-structure. Note that this mesh cannot be converted into a doubly periodic cell transitive Honeycomb structure, as its boundaries do not align. }
\label{fig_teaser1a}
\end{figure}



\subsection{Conversion of Any Mesh into a Link}


The core of our method involves converting any given mesh into a link. Since this conversion is relatively straightforward from a computer graphics perspective, we omit the geometric details from the main paper. Readers can refer to the supplementary document for a full description of the link generation process.

Without loss of generality, we convert each face of the mesh into \ergun{a generalized torus—that is, a 3D-thickened version of the corresponding \textit{face-cycle}. The face-cycles are planar curves constructed as quadratic B-spline curves, computed using the corner positions of each face as control points. This ensures} that face-cycles sharing an edge are tangent to each other at least at the midpoint of that edge. Additionally, we allow neighboring face-cycles to share a straight segment along the common edge, resulting in tangency over an extended region. This promotes smooth visual continuity and facilitates bundling. \ergun{The final toroidal surface, or \textit{toroidal face-cycle}, is generated by sweeping a circular cross-section along an offset version of the original planar curve, such that adjacent toroids are tangent to each other.}

Due to the construction of face-cycles, bundling is achieved using \ergun{toroidal connectors, which are 3D-thickened versions of corresponding \textit{edge-cycles}. 
% These edge-cycles are circles that are centered at the midpoints of selected (labeled) mesh edges and lie in planes perpendicular to the corresponding edge directions. 
These edge-cycles are circles that are centered at the midpoints of selected (labeled) mesh edges and lie in planes perpendicular to the directions of the corresponding edges.
\textit{toroidal edge-cycles}, are generated by sweeping a circular cross-section along the circular edge-cycles.}
For three-dimensional realization, each \ergun{toroidal} edge-cycle can always be chosen large enough to fully enclose the associated face-cycles, as illustrated in Figures~\ref{fig_teaser1a}, \ref{fig_teaser1b}, and \ref{fig_teaser1c}.



\begin{figure}[htpb]
    \centering    
        \begin{subfigure}[t]{0.31\linewidth}
    \includegraphics[width=0.99\linewidth]{images/link_algo/3manwb}
            \caption{\it The original mesh.}
        \label{fig_3D_Mesh1_0}
    \end{subfigure}
        \hfill
       \begin{subfigure}[t]{0.31\linewidth}
    \includegraphics[width=0.99\linewidth]{images/link_algo/3manwb_linked}
            \caption{\it Related fully-linked structure.}
        \label{fig_3D_Mesh1_1}
    \end{subfigure}
        \hfill
         \begin{subfigure}[t]{0.31\linewidth}
    \includegraphics[width=0.99\linewidth]{images/link_algo/3manwb_linked_partial}
            \caption{\it A \ergun{partially-linked structure with single-link property}.}
        \label{fig_3D_Mesh1_2}
    \end{subfigure}
\caption{\it An example of converting a 3-manifold mesh with boundary into a linked-structure. This particular mesh is generated by decomposing a cube using randomly placed Voronoi sites. Note that it cannot be converted into a triply periodic mesh, i.e. a cell-transitive 3D Honeycomb, as its opposite boundaries do not align.}
\label{fig_teaser1c}
\end{figure}

This method of conversion can be applied to any type of mesh to generate a linked-structure. However, our primary interest lies in generating repeating (periodic) structures. It is important to note that repeating links cannot be obtained merely by decomposing space-filling shapes. The two examples shown in Figures~\ref{fig_teaser1a} and \ref{fig_teaser1c} illustrate why this is the case. Although squares and cubes can tile the 2D and 3D spaces, respectively, the decomposed versions of these meshes have mismatched boundaries. As a result, the corresponding linked-structures cannot be extended to fill space periodically. This observation motivates the need for a theoretical framework capable of producing meshes with consistent, repeatable boundary conditions suitable for periodic link generation.
 


\subsection{Bravais Lattices and Wigner-Seitz Cells}

Instead of general meshes, we focus on converting triply periodic 3-manifold meshes into links. The advantage of this approach is that the resulting links will also exhibit triply periodic patterns. To construct these meshes, we employ \textit{Bravais lattices}. In geometry and crystallography, Bravais lattices are defined as infinite $N$-dimensional arrays of discrete points generated by a set of discrete translation operations in $N$-dimensional space, given by:
\[
\mathbf{p}_{\mathbf{n}} = \sum_{i=0}^{N-1} n_i \vec{v}_i
\]
where $\mathbf{n} = (n_0, n_1, \dots, n_i, \dots, n_{N-1}) \in \mathbb{Z}^N$ is any set of $N$ integers, and $\vec{v}_i \in \mathbb{R}^N$ are $N$ linearly independent vectors (not necessarily mutually perpendicular) that span the lattice. 

Note that the lattice points correspond to the vertices of an $N$-dimensional honeycomb formed by tiling $N$-dimensional parallelepipeds. An exciting property of Bravais lattices is that if we use the lattice points as Voronoi sites we obtain a variety of combinatorially distinct classes of $N$-dimensional space-filling tiles, which are called the Wigner–Seitz cells. For instance, in 2D we can obtain both rectangles and hexagons as Wigner–Seitz cells as shown in Figure~\ref{fig_2DWS}. Wigner-Seitz cells in 3D are more interesting. \ergun{There are only five combinatorially distinct classes of polyhedra, all of which are parallelohedra.}: as cube, hexagonal prism, truncated octahedron, elongated dodecahedron, and rhombic dodecahedron (see Figure~\ref{fig_teaser} linked versions of these Wigner-Seitz cells). 		
Note that Wigner-Seitz cells are triply periodic structures. The portions on the Wigner-Seitz cells that are inside of $N$-dimensional parallelepipeds repeat in all directions that are given by  $\vec{v}_i \in \mathbb{R}^N$ (See  Figure~\ref{fig_2DWS} for 2D example). Our theoretical framework to obtain triply periodic meshes is based on this property of Bravais Lattices. 



\subsection{Triply Periodic 3-manifold Meshes with Planar Faces}
\label{sec_TP3MMwPF}

Now, consider the parametric $N$-dimensional parallelepiped defined as:
\[
\mathbf{p}(\mathbf{u}) = \sum_{i=0}^{N-1} u_i \vec{v}_i
\]
where $\mathbf{u} \in [0,1]^N$ is the parameter vector, with each $u_i \in [0,1]$. We can use $\mathbf{u}$ as a local coordinate system. Since this generalizes the idea of 2D texture coordinates, we refer to $\mathbf{u}$ as \textit{$N$-dimensional texture coordinates}.
To generate triply periodic random structures, consider a set of 3D texture coordinates $\mathbf{u}_j \in [0,1]^3$, produced by some method (the specific method is not essential here). Using these texture coordinates, we can compute corresponding positions in 3D space as $\mathbf{p}(\mathbf{u}_j)$. Within the Bravais lattice framework, this yields repeating points in all three directions: $
\mathbf{p}(\mathbf{u}_j) + \mathbf{p}_{\mathbf{n}}
$
for all integer triples $\mathbf{n}$.

Triply periodic planar meshes can be constructed by treating these points as Voronoi sites. Note that computing a Voronoi decomposition requires the points inside a single 3D parallelepiped \ergun{and those in all of its neighboring cells.}

If the $N$-dimensional texture coordinates are generated randomly, the result is a triply periodic random mesh and it can be converted to triply periodic links (Figure~\ref{fig_randommesh}). Moreover, if these coordinates are generated procedurally, some form of visible structure or pattern may also emerge. Both approaches can be useful. However, it is difficult to procedurally label edges when the Voronoi decomposition is derived from randomly or procedurally generated points.
Therefore, in this paper, we focus on the simplest case: \textit{cubical arrays}, which allows more predictable and controllable labeling strategies.




\begin{figure}[htpb]
    \centering    
        \begin{subfigure}[t]{0.49\linewidth}
    \includegraphics[width=0.99\linewidth]{images/EXAMPLES/RANDOM_MESH.jpg}
            \caption{\it The original mesh.}
        \label{fig_randommesh_0}
    \end{subfigure}
    \begin{subfigure}[t]{0.49\linewidth}
    \includegraphics[width=0.99\linewidth]{images/EXAMPLES/RANDOM_MESH_REP.jpg}
            \caption{\it Repeating mesh structure.}
        \label{fig_randommesh_0_REP}
    \end{subfigure}
       \begin{subfigure}[t]{0.49\linewidth}
    \includegraphics[width=0.99\linewidth]{images/EXAMPLES/RANDOM_MESH_FULLY_LINKED.jpg}
            \caption{\it Related fully-linked structure.}
        \label{fig_randommesh_1}
    \end{subfigure}
        \hfill
         \begin{subfigure}[t]{0.49\linewidth}
    \includegraphics[width=0.99\linewidth]{images/EXAMPLES/RANDOM_MESH_PARTIALLY_LINKED.jpg}
            \caption{\it A \ergun{partially-linked structure with single-link property}.}
        \label{fig_randommesh_2}
    \end{subfigure}
\caption{\it An example of linked-structures generated from triply periodic random 3-manifold meshes. This particular mesh is constructed by placing Voronoi sites randomly within a parallelepiped-shaped cell of a transitive 3D honeycomb, which tiles space through translations of the same parallelepiped. Because the opposite boundaries of the mesh align, it produces a triply periodic linked-structure.}
\label{fig_randommesh}
\end{figure}

\subsection{Volumetric Arrays}

Assume that a cubical array of $M \times M \times M$ Bravais points is defined as follows:
\[
\mathbf{p}_\mathbf{m} = \sum_{i=0}^{2} \frac{m_i}{M} \vec{v}_i
\]
where $\mathbf{m} = (m_0,m_1,m_2) \in \{ 0,1,2,\ldots,M-1\}^3$.

Note that each of these points generates a scaled version of the original Wigner-Seitz cell by a factor of $M$, since the new structure corresponds to $M$-scaled versions of the original Bravais vectors. In other words, if these points are used as Voronoi sites, they form geometrically deformed versions of one of the five topologically distinct Wigner-Seitz cells. That is, each point corresponds to a unique cell, and all cells are identical in shape and configuration.

To construct links, we must label the edges of the cells. It is crucial to develop a scheme that avoids redundancy. We observe that each cell can be uniquely assigned labeling information using a single number. We investigated whether it is possible to assign each cell responsibility for labeling a subset of its edges such that all edges in the structure are labeled without overlap. It turns out that this is indeed possible for each of the five combinatorially distinct Wigner-Seitz cells.

Using a brute-force algorithm, we determined the number of edges that need to be labeled for each Wigner-Seitz cell type:

\noindent $\bullet$ \textbf{cP (Cube)}: 3 edges. Each edge connects 4 cells.

\noindent $\bullet$ \textbf{hP (Hexagonal Prism)}: 5 edges. 3 edges connect 4 cells, and 2 edges connect 3 cells.

\noindent $\bullet$ \textbf{cF (Rhombic Dodecahedron)}: 8 edges. Each edge connects 3 cells.

\noindent $\bullet$ \textbf{oF (Elongated Dodecahedron)}: 9 edges. 1 edge connects 4 cells, and 8 edges connect 3 cells.

\noindent $\bullet$ \textbf{hR (Truncated Octahedron)}: 12 edges. Each edge connects 3 cells.

Each edge can be labeled as either 0 or 1, representing the absence or presence of an edge-cycle, respectively. Hence, if $E$ is the number of edges that can be redundantly labeled, each Wigner-Seitz cell can be assigned a number between $0$ and $2^E - 1$, indicating which edges are bundled with edge-cycles. Note that each edge label is independent of the other; therefore, this forms a vector of $E$ length where each element comes from the Finite Field (or Galois Field) $GF(2)$ \cite{lidl1997finite}. In other words, the range is the $GF(2)^E$. Moreover, regardless of shape, all Wigner-Seitz cells are organized in a cubical array uniquely identified by indices \ergun{ $(m_0, m_1, m_2) \in \mathbb{Z}_M^3$, where $\mathbb{Z}_M$ or more formally $\mathbb{Z}/M\mathbb{Z}$ is the ring of integers modulo $M$. This is needed since this structure is cell-transitive. In conclusion, the entire structure can be represented as a single function:
\[
F: \mathbb{Z}_M^3 \rightarrow GF(2)^E.
\]
}
In 3D, the function $F(m_0, m_1, m_2)$ uniquely defines a triply periodic linked-structure. This triply periodic mesh with labeled edges can then be converted into a triply periodic link by turning every face into a \ergun{toroidal} face-cycle and every labeled edge into a \ergun{toroidal} edge-cycle.

We observed that the problem of pattern generation for linked-structures can be significantly simplified by adapting strategies originally developed for woven patterns, owing to the inherent regularity and symmetry of Bravais lattices. In the following section, we discuss rule-based strategies for designing patterns that achieve the single-link property.


\section{Pattern Generation Using Labeling Strategies}

A key advantage of our framework is its ability to generate a wide variety of linked-structures through edge-labeling strategies applied to volumetric arrays of Wigner–Seitz cells. As introduced in Section~\ref{sec_preliminaries}, each cell in the cubical array is assigned a label that encodes which of its $E$ edges are bundled using edge-cycles. These binary labels can be interpreted as an $E$-dimensional vector with elements of the Galois field $GF(2)$, i.e. $GF(2)^E$, and the labeling function 
$F(m_0, m_1, m_2)$
fully specifies the topology of the resulting periodic structure. Note that $GF(2)^E$ is a vector, meaning this function comprises $E$ component functions. We can, therefore, express $F$ as a collection of labeling functions, each responsible for defining edge-cycles in one type of edge:

$$F(m_0, m_1, m_2) = \left(F_0, F_1, \ldots, F_{E-1} \right)$$
where each $F_i$ is a function from $GF(N)^3$ to $GF(2)$. We now proceed to define rules for constructing these labeling functions.

\subsection{Rule-Based Labeling}

In woven structures, the most common rule-based operators are shift operations~\cite{grunbaum1980satins, grunbaum1985catalogue, grunbaum1986extension, grunbaum1988isonemal}. These were originally developed for planar (2D) patterns. In 2D, a shift operation is typically represented as $(N,s)$, where $N$ is a binary number of length $M$, with $0$ denoting a down-cross and $1$ denoting an up-cross. It is important to note that the shift value $s$ cannot be chosen independently of $M$; it must be one of the divisors of $M$ to ensure periodicity.

The shift operator can generate all well-known fabric patterns in 2D. For instance, setting $s=1$ yields all twill patterns; $N=10$ corresponds to plain weaving~\cite{akleman2009cyclic}, and $N=1100$ yields a common twill pattern~\cite{akleman2011cyclic}. However, only twill patterns that can have more ones than zeros can produce single-links (Figure~\ref{fig_4x4x4twill}).

In satin weaves, $N$ is a binary number of the form of either $00\ldots001$ or $11\ldots110$, and $s$ must be relatively prime to $M-1$. Satin patterns where $N=00\ldots001$ do not provide single-links; on the other hand, $N$ patterns where $N=11\ldots110$ guarantee to provide single-links. In weaving, these two are equivalent, but in links, they are significantly different (Figure~\ref{fig_satin}). 

This concept can be extended to 3D volumetric structures by introducing additional shift parameters, resulting in a triplet $(N, s_0, s_1)$, where $N$ remains a binary string $a_0a_1\ldots a_{E-1}$ with each $a_i$ in $\{0,1\}$~\cite{yildiz2024volumetric}. Again, the shift values $s_0$ and $s_1$ must be divisors of $M$ to preserve periodicity.

\begin{figure}[htpb]
    \centering    
    %     \begin{subfigure}[t]{0.31\linewidth}
    % \includegraphics[width=0.99\linewidth]{images/EXAMPLES/1000_1_1.jpg}
    %         \caption{\it The twill pattern $(1000,1,1)$ in two directions produces multi-links.}
    %     \label{fig_4x4x4twill_0}
    % \end{subfigure}
    %     \hfill
       \begin{subfigure}[t]{0.49\linewidth}
    \includegraphics[width=0.99\linewidth]{images/EXAMPLES/1100_1_1.jpg}
            \caption{\it The twill pattern $(1100,1,1)$ in two directions produces multi-link.}
        \label{fig_4x4x4twill_1}
    \end{subfigure}
        \hfill
         \begin{subfigure}[t]{0.49\linewidth}
    \includegraphics[width=0.99\linewidth]{images/EXAMPLES/1110_1_1.jpg}
            \caption{\it The twill pattern $(1110,1,1)$ in two directions produces single-link.}
        \label{fig_4x4x4twill_2}
    \end{subfigure}
\caption{\it Two examples of twill patterns for $M=4$. Note that although both weaving patterns are fabrics and hang-together, only $(1110,1,1)$ can produce a single-link.The equivalent fabric pattern $(0001,1,1)$ also does not produce a single-link. In other words, unlike fabrics in links the number of $1$'s has to be more than the number of $0$'s. }
\label{fig_4x4x4twill}
\end{figure}

In our framework, since the label is a vector of binary numbers, we extend this concept to $E$ distinct shift operations by using functions. Let $F_i(x, 0, 0) = \{a_0, a_1, \ldots, a_{M-1}\}$ denote a one-dimensional binary array indexed by $x \in GF(M)$, where each $a_i$ belongs to $GF(2)$. This is analogous to the $N$ string in woven structures. The remaining values of the function are determined by applying a shift operation to this array. This can be achieved by operations over indices. An example is given as follows:
$$F_i(x, y, z) = \{a_{ys_y + zs_z}, a_{1 + ys_y + zs_z}, \ldots, a_{M-1 + ys_y + zs_z}\}.$$
Here, addition is taken modulo $M$ since we are operating in $GF(M)$, and $s_y$ and $s_z$ are the amounts of shift in the $y$ and $z$ directions, respectively.

\begin{figure}[htpb]
    \centering    
    %     \begin{subfigure}[t]{0.31\linewidth}
    % \includegraphics[width=0.99\linewidth]{images/EXAMPLES/00000001_3_3.jpg}
    %         \caption{\it The satin pattern $(00000001,3,3)$ in two directions produces multi-links.}
    %     \label{fig_satin_0}
    % \end{subfigure}
    %     \hfill
       \begin{subfigure}[t]{0.49\linewidth}
    \includegraphics[width=0.99\linewidth]{images/EXAMPLES/11111110_3_3.jpg}
            \caption{\it The satin pattern $(11111110,3,3)$ in two directions produces single-link.}
        \label{fig_satin_1}
    \end{subfigure}
        \hfill
         \begin{subfigure}[t]{0.49\linewidth}
    \includegraphics[width=0.99\linewidth]{images/EXAMPLES/11111110_3_5.jpg}
            \caption{\it The satin pattern $(11111110,3,5)$ in two directions produces single-link.}
        \label{fig_satin_2}
    \end{subfigure}
\caption{\it Examples of satin patterns that can produce single-links for $M=8$. Note that although $(00000001,3,3)$ and $(00000001,3,5)$ are equivalent as woven patterns, but they cannot produce single-links. In other words, although the naming the patterns comes from weaving, this is a different problem.  }
\label{fig_satin}
\end{figure}



In woven structures, another type of operation is the alternating shifts followed by inversion of the binary values. Inversions can be obtained functionally using operations over the values. An example is given as follows:
$$F_i(x, y, z) = \{x + y + a_{ys_y + zs_z}, \ldots, x + y + a_{M-1 + ys_y + zs_z}\}.$$
Note that all operations over values must be carried out in $GF(2)$ since the values are elements of $GF(2)$.

\subsection{Random Labeling}

To explore less regular but still cohesive forms, one can define $F$ as a stochastic function. For instance, a 3D Perlin noise or cellular noise function thresholded at a given level can determine edge-cycle labeling.
While such strategies sacrifice strict regularity, they can create richly varied linked-structures useful for mimicking organic or naturally irregular forms. Figures~\ref{fig_randomHp}, and ~\ref{fig_randomCf} show some examples of random labeling. 



\subsection{Single-Link Structures}

We define a link as a \textit{single-link structure} if every face-cycle is connected, either directly or indirectly, to at least one other face-cycle through a sequence of edge-cycles, such that the union of all face-cycles and edge-cycles forms a single connected component in $\mathbb{R}^3$.

Formally, this system can be represented as a hypergraph $\mathbb{G} = (\mathbb{V}, \mathbb{E})$, where $\mathbb{V}$ is the set of all faces of the original mesh, corresponding to face-cycles, and $\mathbb{E}$ is the set of all labeled edges of the original mesh, consisting of all subsets of $\mathbb{V}$ that share a labeled edge. In other words, $\mathbb{E}$ represents the edge-cycles. The link is a single-link structure if the hypergraph $\mathbb{G}$ is connected. According to this definition, fully-linked structures naturally become single-linked if the corresponding mesh is connected. In contrast, partially-linked structures do not necessarily guarantee a single-link outcome.
Figures~\ref{fig_randomHp}, and ~\ref{fig_randomCf} also show examples of single links. 


% \begin{figure}[htpb!]
%     \centering        
%         \begin{subfigure}[t]{0.49\linewidth}
%     \includegraphics[width=0.99\linewidth]{images/EXAMPLES/CP_RANDOM_0.6}
    
%             \caption{\it Random-label single links that use only 60\% of edge-cycles.}
%         \label{fig_randomCp_0}
%     \end{subfigure}
%         \hfill
%        \begin{subfigure}[t]{0.49\linewidth}
%     \includegraphics[width=0.99\linewidth]{images/EXAMPLES/CP_RANDOM_0.75}
%             \caption{\it Random-label single links that use 75\% of edge-cycles.}
%         \label{fig_randomCp_1}
%     \end{subfigure}
%     %     \hfill
%     %      \begin{subfigure}[t]{0.49\linewidth}
%     % \includegraphics[width=0.99\linewidth]{images/EXAMPLES/CP_RANDOM_0.9}
%     %         \caption{\it Random-label single links that use 90\% of edge-cycles.}
%     %     \label{fig_randomCp_2}
%     % \end{subfigure}
% \caption{\it Examples of random labeled single links for Cp case. }
% \label{fig_randomCp}
% \end{figure}

\subsection{Ensuring Single Links}
\label{sec_ESL}

Not all labeling functions $F(m_0, m_1, m_2)$ guarantee the formation of a single linked-structure. Successful configurations require verification of topological connectivity through hypergraph analysis. The verification process involves constructing a hypergraph $G = (V, E)$, where the vertices $V$ represent face-cycles and the hyperedges $E$ correspond to edge-cycles. A breadth-first search is then performed on $G$ to detect disconnected components. The labeling function $F$ is iteratively modified by strategically adding edge-cycle labels until $G$ becomes fully connected. 

For random labeling, the process begins with a sparse distribution of edge-cycles, identifies disconnected components through hypergraph traversal, and iteratively adds critical edge-cycles at the boundaries of these components. This strategy converges quickly, typically requiring $O(M)$ modifications for $M^3$ systems. Rule-based labeling allows for limited rule violations at domain boundaries to ensure global connectivity, analogous to single-cycle weaving strategies~\cite{xing2010single}. The minimum percentage of labeled edges required to achieve a single link depends on the specific structure. In most cases, single-link configurations were achieved with 60\% of the edges labeled; however, this threshold could not be met in some instances (See Figures~\ref{fig_randomHp}, and~\ref{fig_randomCf} for examples of three such randomly labeled single-link structures). 



\begin{figure}[htpb!]
    \centering        
        \begin{subfigure}[t]{0.49\linewidth} \includegraphics[width=0.99\linewidth]{images/EXAMPLES/HP_RANDOM_0.6}
            \caption{\it Random-label single links that use only 60\% of edge-cycles.}
        \label{fig_randomHp_0}
    \end{subfigure}
        \hfill
       \begin{subfigure}[t]{0.49\linewidth}
    \includegraphics[width=0.99\linewidth]{images/EXAMPLES/HP_RANDOM_0.75}
            \caption{\it Random-label single links that use 75\% of edge-cycles.}
        \label{fig_randomHp_1}
    \end{subfigure}
    %     \hfill
    %      \begin{subfigure}[t]{0.49\linewidth}
    % \includegraphics[width=0.99\linewidth]{images/EXAMPLES/HP_RANDOM_0.9}
    %         \caption{\it Random-label single links that use 90\% of edge-cycles.}
    %     \label{fig_randomHp_2}
    % \end{subfigure}
\caption{\it Examples of random labeled single links for Hp case. }
\label{fig_randomHp}
\end{figure}

In rule-based labeling, one may prefer not to modify $F$, as doing so would alter the underlying rule and disrupt the intended pattern. For such cases, we analyzed the conditions under which rule-based configurations naturally yield connected structures. We observed that the number of $1$'s in the rule must exceed the number of $0$'s, and the shift value must be relatively prime to $M$ to ensure connectivity along the specified direction. \ergun{Otherwise, the label pattern may form a subgroup that fails to span the entire structure.} Applying this rule in a single direction produces planar single-links within each layer. In the cubic lattice case (Cp), applying compatible rules in two different directions enables the construction of single-links, as shown in Figures~\ref{fig_4x4x4twill} and~\ref{fig_satin}. For other lattice types, additional analysis is required to ensure global connectivity.

% Connectivity enforcement works equally well for both periodic and aperiodic structures, though periodic systems benefit from translational symmetry in label propagation. The verification process scales as $O(N)$ for $N$ face-cycles when using union-find data structures with path compression \cite{tarjan1975efficiency}.

\begin{figure}[htpb!]
    \centering        
       \begin{subfigure}[t]{0.49\linewidth}
    \includegraphics[width=0.99\linewidth]{images/EXAMPLES/CF_RANDOM_0.75}
            \caption{\it Random-label single links that use 75\% of edge-cycles.}
        \label{fig_randomCf_1}
    \end{subfigure}
        \hfill
         \begin{subfigure}[t]{0.49\linewidth}
    \includegraphics[width=0.99\linewidth]{images/EXAMPLES/CF_RANDOM_0.9}
            \caption{\it Random-label single links that use 90\% of edge-cycles.}
        \label{fig_randomCf_2}
    \end{subfigure}
\caption{\it Examples of randomly labeled single links for the hR case. For hR, we were unable to achieve a single link using only 60\% of the edge-cycles.}
\label{fig_randomCf}
\end{figure}


\section{Implementation and Results}

We developed an integrated Blender plugin that combines mesh-to-link conversion with Voronoi decomposition capabilities for lattice structures. The tool enables direct transformation of arbitrary input meshes into linked-structures while maintaining an interactive workflow for Bravais lattice manipulation. Users can dynamically deform lattice parameters and observe real-time updates to corresponding Wigner-Seitz cells, which are subsequently converted into desired link configurations through customizable edge-labeling functions.

The plugin's core functionality operates through four synchronized components: 1) a lattice deformation module handling Bravais structure manipulation, 2) our Voronoi plug-in that is used for generating Wigner-Seitz cells, 3) a labeling engine applying user-specified edge-bundling patterns, and 4) connectivity validation to obtain single-links. The connectivity validation occurs during generation using the union-find algorithm detailed in Section~\ref{sec_ESL}.

All featured structures in this paper were generated and rendered within the same software environment.
Our software demonstrates significant versatility in link generation, as evidenced by the diverse structures shown in the examples provided in this paper, such as Figure~\ref{fig_teaser}. These examples span multiple Bravais lattice types and labeling strategies, all produced through direct mesh conversion within the Blender environment. The tight integration between modeling and visualization subsystems ensures geometric fidelity between computational representations and final rendered outputs.


\section{Conclusion, Discussion and Future Work}


\textit{Overview:} The proposed framework for constructing triply periodic linked-structures introduces a novel approach to generating geometrically simple yet topologically robust links from triply periodic meshes. By treating polygonal mesh geometry and topology as design parameters, the method achieves a balance between flexibility and simplicity. The transformation of mesh faces into unknotted face-cycles, combined with edge-cycles as bundling elements, offers an intuitive and powerful mechanism for designing linked-structures.

\textit{Physical Evaluations:} This labeling-based formulation defines a high-dimensional design space for volumetric linked-structures. With $2^E$ possible configurations per cell and $M^3$ cells, even modest volumes yield combinatorially rich pattern sets. This design space can be explored algorithmically to identify linked-structures with desirable properties. We are currently collaborating with an engineering group interested in these structures. This group has developed a system for evaluating certain physical properties computationally.

\textit{Physical Constructions:} Evaluating physical models will also be an important next step. There are three general methods for constructing these structures in practice: (1) 3D printing the entire structure as a single piece; (2) 3D printing individual face-cycles and bundling them using commercially available materials such as spiral wraps; and (3) assembling them entirely from off-the-shelf components such as metal rings, cable ties, and spiral wraps. Among these, method (3) becomes increasingly difficult for complex structures. Method (2) is relatively more feasible, though it still demands substantial manual effort.

\textit{Other Labeling Strategies:} The group structure of $GF(M)^3$ supports more sophisticated labeling strategies involving algebraic operations such as reflection, rotation, translation, and tensor contraction. For instance, tensorial combinations of bit patterns across neighboring cells can produce structured yet nontrivial edge-cycle arrangements, ideal for generating visually intricate but topologically coherent linked-structures. These algebraic strategies are particularly effective at leveraging the symmetries inherent in Bravais lattices.

\textit{Topology Limitations:} The current method is capable of generating only a restricted class of links and cannot represent more complex configurations such as the Borromean rings. This limitation is intentional: although more general techniques for constructing links exist, they often result in geometrically complex shapes that are less suitable for engineering applications. Our framework prioritizes simplicity and realizability.

\begin{figure}[htpb]
    \centering    
        \begin{subfigure}[t]{0.485\linewidth}
    \includegraphics[width=0.99\linewidth]{images/EXAMPLES/HIERARCHICAL_BASE.jpg}
            \caption{\it 2 disconnected but linked meshes.}
        \label{fig_hierarchy_0}
    \end{subfigure}
        \hfill
       \begin{subfigure}[t]{0.485\linewidth}
    \includegraphics[width=0.99\linewidth]{images/EXAMPLES/HIERARCHICAL_LINKS.jpg}
            \caption{\it 2 links that are linked together.}
        \label{fig_hierarchy_1}
    \end{subfigure}
\caption{\it An example of hierarchical links. }
\label{fig_hierarchy}
\end{figure}


\textit{Biomimetic Hierarchical Linking.} \ergun{Figure~\ref{fig_hierarchy} demonstrates an exciting direction for future research is to obtain geometric interlocking using hierarchical, multiscale architectures inspired by biological systems. By implementing nested polycatanated linking across scales, it may be possible to achieve geometric interlocking with multiple links.} These can provide enhanced mechanical properties such as fracture resistance and crack deflection, analogous to mechanisms observed in natural materials like nacre and collagen networks~\cite{Chen2024deployable}.


\bibliographystyle{cag-num-names}
\bibliography{references,references2}

\end{document}

